\documentclass[12pt]{article} % Starts an article
\usepackage{amsmath} % Imports amsmath
\title{Andy Atkeson} % Title

\begin{document} % Begins a document
  \maketitle
  Andrew Atkeson (1991) designed a model to explain how, in defiance of the
  pattern predicted by complete markets models, low output realizations in various countries in the mid-1980s prompted international lenders to ask those
  countries for net repayments. A complete markets model would have net flows
  to a borrower during periods of bad endowment shocks. Atkeson’s idea was
  that information and enforcement problems could produce the observed outcome. Thus, Atkeson’s model combines two features of the models we have seen
  in chapter 22: incentive problems from private information and participation
  constraints coming from enforcement problems.
  Atkeson showed that the optimal contract handles enforcement and information problems through the shape of the repayment schedule, thereby indirectly
  manipulating continuation values. Continuation values respond only by updating  a single state variable, a measure of resources available to the borrower, that
  appears in the optimum value function, which in turn is affected only through
  the repayment schedule. Once this state variable is taken into account, promised
  values do not appear as independently manipulated state variables. 
  Atkeson’s model brings together several features. He studies a “borrower”
  who by himself is situated like a planner in a stochastic growth model, with the
  only vehicle for saving being a stochastic investment technology. Atkeson adds
  the possibility that the planner can also borrow subject to both participation
  and information constraints.

 
\end{document}